\documentclass[../../main.tex]{subfiles}% 注意这里的文件路径不能用 ./main.tex ,否则用latexmk编译子文件会报错
\graphicspath{{\subfix{./image/}}{\subfix{../../image/}}} % 指定图片引用路径,后续可以直接使用图片文件名
% 如果图片存放在上级目录,则使用 ../../image/ 作为备用路径

% 例如:
% \begin{figure}[H]
% \centering
% \includegraphics[width=0.3\textwidth]{图.png}
% \caption{}
% \label{figure:图}
% \end{figure}
% 注意:上述\label{}一定要放在\caption{}之后,否则引用图片序号会只会显示??.

\begin{document}

\section{子群与商群}

\begin{definition}
设 \( A, B \) 是群 \( G \) 的两个子集, 约定
\begin{align*}
AB = \{ab|a \in A, b \in B\} , A^{-1} = \{a^{-1}|a \in A\}.
\end{align*}
特别地, 当 \( A = \{a\} \) 为单点集时, 记 \( AB = aB \), \( BA = Ba \). 当然这些符号对半群与幺半群可同样使用.
\end{definition}

\begin{proposition}\label{proposition:有限群乘积的阶}
设$G$是群,$H$是$G$的有限子集,$a,b\in G$,则
\begin{align*}
|H|=|aH| = |Ha| = |aHb| .
\end{align*}
\end{proposition}
\begin{proof}
令
\begin{align*}
\varphi : &H \longrightarrow aH, \\
&h \longmapsto ah,\quad \forall h \in H.
\end{align*}
容易验证$\varphi$是$H$到$aH$的双射.故$|H|=|aH|$.其他同理可证.
\end{proof}

\begin{definition}
群 \( G \) 的非空子集 \( H \) 若对 \( G \) 的运算也构成一个群, 则称为 \( G \) 的\textbf{子群},记作$H\leqslant G$.若还有$H\neq G$,则称$H$为$G$的\textbf{真子群},记作$H<G$.
\end{definition}
\begin{remark}
显然, \( H = \{1\} \)(1 为 \( G \) 的幺元) 与 \( H = G \) 均为 \( G \) 的子群, 称为 \( G \) 的\textbf{平凡子群}, 其他的子群称为\textbf{非平凡子群}.
\end{remark}

\begin{definition}
若半群 \( S \) 的非空子集 \( S_1 \) 对 \( S \) 的运算也是半群, 则称 \( S_1 \) 为 \( S \) 的\textbf{子半群}.

若幺半群 \( M \) 的子集 \( Q \) 对 \( M \) 的运算也是幺半群且 \( M \) 的幺元 \( 1 \in Q \), 则称 \( Q \) 为 \( M \) 的\textbf{子幺半群}.
\end{definition}

\begin{theorem}\label{theorem:子群的充要条件--抽象代数}
设 \( H \) 是群 \( G \) 的非空子集, 则下列条件等价:
\begin{enumerate}[(1)]
\item\label{theorem:子群的充要条件--抽象代数-1} \( H \) 是 \( G \) 的子群.

\item\label{theorem:子群的充要条件--抽象代数-2} \( 1 \in H \); 若\( a \in H \), 则 \( a^{-1} \in H \); 若\( a, b \in H \), 则 \( ab \in H \).

\item\label{theorem:子群的充要条件--抽象代数-3} 对\(\forall a, b \in H \), 有 \( ab \in H \), \( a^{-1} \in H \).

\item\label{theorem:子群的充要条件--抽象代数-4} 对\( a, b \in H \), 有 \( ab^{-1} \in H \).

\item\label{theorem:子群的充要条件--抽象代数-5} 对\( a, b \in H \), 有 \( a^{-1}b \in H \).
\end{enumerate}
\end{theorem}
\begin{proof}
\ref{theorem:子群的充要条件--抽象代数-1} $\Rightarrow$ \ref{theorem:子群的充要条件--抽象代数-2}. 由 \( H \) 对 \( G \) 的乘法构成群知 \( a, b \in H \), 则 \( ab \in H \). 又 \( H \) 有幺元 \( 1' \), 即有 \( 1' \cdot 1' = 1' \). 设 \( 1' \) 在 \( G \) 中的逆元为 \( 1'^{-1} \), 则有
\[
1 = 1' \cdot 1'^{-1} = (1' \cdot 1') \cdot 1'^{-1} = 1',
\]
故 \( 1 \in H \). 设 \( a \) 在 \( H \) 中的逆元为 \( a' \), 于是 \( aa' = 1' = 1 \), 即 \( a' = a^{-1} \), 故 \( a^{-1} \in H \). 由此知 \( (2) \) 成立, 而且 \( H \) 的幺元是 \( G \) 的幺元. \( a \in H \), \( a \) 在 \( H \) 中的逆元与在 \( G \) 中的逆元一致.

\ref{theorem:子群的充要条件--抽象代数-2} $\Rightarrow$ \ref{theorem:子群的充要条件--抽象代数-3}. 这是显然的.

\ref{theorem:子群的充要条件--抽象代数-3} $\Rightarrow$ \ref{theorem:子群的充要条件--抽象代数-4}. 若 \( a, b \in H \), 故 \( a, b^{-1} \in H \), 故 \( ab^{-1},a^{-1}b \in H \).

\ref{theorem:子群的充要条件--抽象代数-4} $\Rightarrow$ \ref{theorem:子群的充要条件--抽象代数-1} . 由 \( H \neq \varnothing \) 知 \( \exists a \in H \), 因而 \( 1 = aa^{-1} \in H \). 又由 \( 1, a \in H \) 知 \( a^{-1} = 1 \cdot a^{-1} \in H \). 又若 \( a, b \in H \), 由 \( b^{-1} \in H \) 得 \( ab = a(b^{-1})^{-1} \in H \). 由此可知 \( G \) 的乘法也是 \( H \) 的乘法. 对 \( H \) 而言有幺元 \( 1 \); 对 \( a \in H \) 有逆元 \( a^{-1} \); 结合律显然成立. 故 \( H \) 是 \( G \) 的子群.

\ref{theorem:子群的充要条件--抽象代数-5} $\Rightarrow$ \ref{theorem:子群的充要条件--抽象代数-1}. 由 \( H \neq \varnothing \) 知 \( \exists a \in H \), 因而 \( 1 = aa^{-1} \in H \). 又由 \( 1, a \in H \) 知 \( a^{-1} = 1 \cdot a^{-1} \in H \). 又若 \( a, b \in H \), 由 \( a^{-1} \in H \) 得 \( ab = (a^{-1})^{-1}b \in H \). 由此可知 \( G \) 的乘法也是 \( H \) 的乘法. 对 \( H \) 而言有幺元 \( 1 \); 对 \( a \in H \) 有逆元 \( a^{-1} \); 结合律显然成立. 故 \( H \) 是 \( G \) 的子群.
\end{proof}

\begin{theorem}\label{theorem:子群陪集的基本性质}
设$A,B,C,H$是群$G$的非空子集,$g,a,b\in G$,则
\begin{enumerate}[(1)]
\item\label{theorem:子群陪集的基本性质-1} $A(BC)=(AB)C$.

\item\label{theorem:子群陪集的基本性质-2} $gA=gB\text{或}Ag=Bg\iff A=B$.

\item\label{theorem:子群陪集的基本性质-3} \( H \) 是 \( G \) 的子群$\iff HH = H , H^{-1} = H\iff H^{-1}H = H$.

\item\label{theorem:子群陪集的基本性质-4} $ab\in H\iff a\in Hb^{-1}\iff b\in a^{-1}H$.

\item $x\in gH \iff x^{-1}\in H^{-1}g^{-1}$,$\,\,x\in Hg \iff x^{-1}\in g^{-1}H^{-1}$,$\,\,x\in HgK\iff x^{-1}\in K^{-1}g^{-1}H^{-1}$.
\end{enumerate}
\end{theorem}
\begin{proof}
\begin{enumerate}[(1)]
\item 

\item 

\item 

\item 只需注意到$ab\in H\Longleftrightarrow ab=h\left( h\in H \right) \Longleftrightarrow a=hb^{-1}\in Hb^{-1}\Longleftrightarrow b\in a^{-1}H$.

\item 
\end{enumerate}
\end{proof}

\begin{proposition}\label{proposition:抽象代数---子群的基本性质}
\begin{enumerate}[(1)]
\item\label{proposition:抽象代数---子群的基本性质-0} 设$K\leqslant H\leqslant G$,则$K\leqslant G$.

\item\label{proposition:抽象代数---子群的基本性质-1} 设$H$是群$G$的非空有限子集,则$H$是$G$的子群的充分必要条件是$H$关于$G$的运算封闭.

\item\label{proposition:抽象代数---子群的基本性质-2} 若$G$是一个群,则$G$的任意子群的交$\bigcap_{H<G}{H}$也是$G$的子群.

\item\label{proposition:抽象代数---子群的基本性质-3} 若$H_1,H_2$都是群$G$的子群且$H_2\subseteq H_1$,则$H_2$也是$H_1$的子群.

\item\label{proposition:抽象代数---子群的基本性质-4} 设$H,K$是群$G$的两个子群. 则$H\cup K$是$G$的子群的充要条件是$H\subseteq K$或$K\subseteq H$. 并且群$G$不能被它的两个真子群所覆盖.

\item\label{proposition:抽象代数---子群的基本性质-5} 设$A$和$B$是有限群$G$的两个非空子集. 若$|A|+|B|>|G|$, 则$G=AB$.
\end{enumerate}
\end{proposition}
\begin{remark}
在这个\rrefpro{proposition:抽象代数---子群的基本性质}{proposition:抽象代数---子群的基本性质-4}中,群$G$可能被它的三个真子群所覆盖. 例如, 群
\begin{align*}
U(8)=\{\overline{1},\overline{3},\overline{5},\overline{7}\}.
\end{align*}
易知
\begin{align*}
H=\{\overline{1},\overline{3}\}, \quad J=\{\overline{1},\overline{5}\}, \quad K=\{\overline{1},\overline{7}\}
\end{align*}
都是$U(8)$的真子群, 且$U(8)=H\cup J\cup K$.
\end{remark}
\begin{proof}
\begin{enumerate}[(1)]
\item 证明是显然的.

\item 必要性显然, 下证充分性.
因为$H$关于$G$的运算封闭, 所以$G$的运算是$H$的代数运算. 又因为$G$的运算满足结合律, 所以$H$的运算也满足结合律.

设$H=\{a_1,a_2,\cdots,a_n\}$. 对任意的$a\in H$, 记$Ha=\{a_1a,a_2a,\cdots,a_na\}$, 则$Ha\subseteq H$.
于是$a_ia\neq a_ja(i\neq j)$.否则,由$a_ia = a_ja(i\neq j)$可得
\begin{align*}
a_i=a_i(aa^{-1})=(a_ia)a^{-1}=(a_ja)a^{-1}=a_j(aa^{-1})=a_j,
\end{align*}
显然矛盾!
由此推出$|Ha|=n=|H|$, 于是$Ha=H$. 这样, 对任意的$a,b\in H$, 因为$Ha=H$, 所以必有$a_i\in H$, 使$a_ia=b$. 这说明, 对任意的$a,b\in H$, 方程$xa=b$在$H$中必有解. 同理可证, 方程$ay=b$在$H$中也有解. 从而, 由\refthe{theorem:群的充要条件--方程有解}知$H$为群.

\item 设$I$为任一(有限或无限的)指标集, $\{H_i\mid i\in I\}$为群$G$的一些子群的集合, 令
\begin{align*}
J = \bigcap\limits_{i\in I} H_i,
\end{align*}
因为$1\in H_i(\forall i\in I)$, 所以$1\in \bigcap\limits_{i\in I} H_i$, 从而$J$非空;
对$\forall a,b\in J$, 有$a,b\in H_i(\forall i\in I)$. 由于$H_i < G$, 因此$ab^{-1}\in H_i(\forall i\in I)$, 于是$ab^{-1}\in J$.
这就证明了$\bigcap\limits_{i\in I} H_i$为$G$的子群.

\item 由$H_2$是$G$的子群知$ab^{-1} \in H_2$，$\forall a,b \in H_2$.又$H_2 \subseteq H_1$，故$H_2$也是$H_1$的子群.

\item 充分性显然, 下证必要性.
设$H\cup K$是$G$的子群. 如果$H\subseteq K$, 则结论成立. 如果$H\nsubseteq K$, 则存在$h\in H$, 使$h\notin K$. 由于$H\cup K$为$G$的子群, 因此对任意的$k\in K$, 有$hk\in H\cup K$. 从而必有$hk\in H$或$hk\in K$. 如果$hk\in K$, 则$h=hk\cdot k^{-1}\in K$, 这与$h$的选取矛盾. 从而必有$hk\in H$, 由此推出$k=h^{-1}\cdot hk\in H$. 由$k$的任意性知$K\subseteq H$. 这就证明了必要性.

设$H,K$是群$G$的两个子群. 如果$G=H\cup K$, 由前面所证, 应有$H\subseteq K$或$K\subseteq H$, 于是有$G=K$或$G=H$. 这说明$H,K$不可能都是$G$的真子群. 因此群$G$不能被它的两个真子群所覆盖.

\item $\forall ab\in AB$,则$a,b\in G$,由$G$是群知$ab\in G$.故$AB\subseteq G$.
$\forall\,g\in G$, $|A^{-1}g|=|A|$. 因$|A|+|B|>|G|$, 故$A^{-1}g\cap B\neq\varnothing$, 于是有$a\in A, b\in B$使得$b=a^{-1}g$, 即$g=ab$. 这表明$G\subseteq AB$.
故$G = AB$.
\end{enumerate}
\end{proof}

\begin{example}
设$A, B$是群$G$的两个子群且$G=AB$. 如果子群$C$包含$A$, 则$C=A(B\cap C)$.
\end{example}
\begin{proof}
设$c\in C$,由$C\subseteq G=AB$知存在$a\in A$,$b\in B$,使得$c=ab$.从而由$C<G$知$b=a^{-1}c\in C$,故$b\in B\cap C$.因此$c=ab\in A\left( B\cap C \right)$.于是$C\subseteq A\left( B\cap C \right)$.
又设$a_1b_1\in A\left( B\cap C \right)$,则$a_1\in A\subseteq C$,$b_1\in C$.由$C<G$知$a_1b_1\in C$.故$C\supseteq A\left( B\cap C \right)$.
综上$C=A\left( B\cap C \right)$.
\end{proof}

\begin{theorem}
\begin{enumerate}
\item 设 \( V \) 是数域 \( \mathbb{P} \) 上的 \( n \) 维线性空间.\( S_V \) 为 \( V \) 上的全变换群, \( GL(V) \) 表示 \( V \) 上所有可逆线性变换的集合, 则 \( GL(V) \) 为 \( S_V \) 的子群, 称为线性空间 \( V \) 的\textbf{一般线性群}. 

又设 \( SL(V) \) 为 \( V \) 上所有行列式等于 1 的线性变换的集合, 则 \( SL(V) \) 是 \( GL(V) \)(同时也是 \( S_V \)) 的子群, 称为\textbf{特殊线性群}.

\item 设 \( V \) 是 \( n \) 维 Euclid 空间. 以 \( O(V) \) 表示 \( V \) 上所有正交变换的集合, \( SO(V) \) 表示所有行列式等于 1 的正交变换的集合, 则 \( O(V) \) 是 \( GL(V) \) 的子群, \( SO(V) \) 是 \( O(V) \) 的子群. \( O(V) \) 称为 \( V \) 的\textbf{正交变换群}, 简称\textbf{正交群}, \( SO(V) \) 称为\textbf{转动群}(或\textbf{特殊正交变换群}、\textbf{特殊正交群}).
\end{enumerate}
\end{theorem}
\begin{remark}
将上述$S_V$换成数域$\mathbb{P}$上的全体方阵构成的乘法群,线性变换换成方阵,结论也成立.
\end{remark}
\begin{proof}

\end{proof}

\begin{definition}
设 \( H,K \) 是群 \( G \) 的子群, 又 \( a \in G \). 集合 \( aH \) 与 \( Ha \) 分别称为以 \( a \) 为代表的 \( H \) 的\textbf{左陪集}与\textbf{右陪集}.集合$HaK$称为以$a$为代表元的$H,K$的\textbf{双陪集}.
\end{definition}

\begin{theorem}\label{theorem:抽象代数-定理 1.3.2}
设 \( H ,K\) 是群 \( G \) 的非空子集.则
\begin{enumerate}[(1)]
\item\label{theorem:抽象代数-定理 1.3.2-1} 由
\[
aRb \iff a^{-1}b \in H
\]
所确定的 \( G \) 中的关系 \( R \) 是一个等价关系的充要条件是$H$是$G$的子群. 并且当$H$是$G$的子群时, \( g \) 所在的等价类为 \( gH \), 故 \( H \)的左陪集族 \( \{gH:g\in G\} \)(集合无相同元素) 是 \( G \) 的一个分划.即
\begin{align*}
G = \bigsqcup_{g\in G}gH,
\end{align*}
其中$g$取遍不同$H$的左陪集的代表元.

\item\label{theorem:抽象代数-定理 1.3.2-2} 由
\[
aRb \iff ab^{-1} \in H
\]
所确定的 \( G \) 中的关系 \( R \) 是一个等价关系的充要条件是$H$是$G$的子群. 并且当$H$是$G$的子群时, \( g \) 所在的等价类为 \( Hg \), 故 \( H \)的右陪集族 \( \{Hg:g\in G\} \)(集合无相同元素) 是 \( G \) 的一个分划.即
\begin{align*}
G = \bigsqcup_{g\in G}Hg,
\end{align*}
其中$g$取遍不同$H$的右陪集的代表元.

\item\label{theorem:抽象代数-定理 1.3.2-3} 由
\[
xRy \iff x\in HyK
\]
所确定的 \( G \) 中的关系 \( R \) 是一个等价关系的充要条件是$H,K$是$G$的子群. 并且当$H,K$是$G$的子群时, \( g \) 所在的等价类为 \( HgK \), 故 \( H \)的右陪集族 \( \{HgK:g\in G\} \)(集合无相同元素) 是 \( G \) 的一个分划.即
\begin{align*}
G = \bigsqcup_{g\in G}HgK,
\end{align*}
其中$g$取遍不同$H,K$的双陪集的代表元.
\end{enumerate}
\end{theorem}
\begin{proof}
\begin{enumerate}[(1)]
\item {\heiti 充分性:}由 \( a^{-1}a \in H \) 知 \( aRa(\forall a \in G) \). 又设 \( aRb \), 即 \( a^{-1}b \in H \), 故 \( (a^{-1}b)^{-1} = b^{-1}a \in H \), 即 \( bRa \). 再设 \( aRb, cRb \), 即 \( a^{-1}b, b^{-1}c \in H \), 故 \( a^{-1}c = (a^{-1}b)(b^{-1}c) \in H \), 即 \( aRc \). 这样知 \( R \) 是等价关系. 又由
\[
gRx \iff \ g^{-1}x \in H \iff \ x=g(g^{-1}x) \in gH
\]
知 \( g \) 所在的等价类为 \( gH \). 由\refthe{Set Theory-theorem:等价类就和集合的分划对应-定理1.1.1}知 \( \{gH:g\in G\} \) 为 \( G \) 的一个分划.

{\heiti 必要性:}由等价关系的定义和\refthe{theorem:子群的充要条件--抽象代数}是显然的.

\item {\heiti 充分性:}由 \( aa^{-1} \in H \) 知 \( aRa(\forall a \in G) \). 又设 \( aRb \), 即 \( ab^{-1} \in H \), 故 \( (ab^{-1})^{-1} = ba^{-1} \in H \), 即 \( bRa \). 再设 \( aRb, cRb \), 即 \( ab^{-1}, cb^{-1} \in H \), 故 \( ac^{-1} = (ab^{-1})(bc^{-1}) \in H \), 即 \( aRc \). 这样知 \( R \) 是等价关系. 又由
\[
xRg \iff \ xg^{-1} \in H \iff \ x=(xg^{-1})g \in Hg
\]
知 \( g \) 所在的等价类为 \( Hg \). 由\refthe{Set Theory-theorem:等价类就和集合的分划对应-定理1.1.1}知 \( \{Hg:g\in G\} \) 为 \( G \) 的一个分划.

{\heiti 必要性:}由等价关系的定义和\refthe{theorem:子群的充要条件--抽象代数}是显然的.

\item {\heiti 充分性:}对 $\forall x\in G$, 由 $x=1\cdot x\cdot 1\in HxK$ 知 $xRx$.
设 $x,y\in G$ 且 $xRy$, 则 $x\in HyK$, 于是存在 $h\in H,k\in K$, 使得 $x=hyk$. 从而
\begin{align*}
y=h^{-1}xk^{-1}\in HxK\Longrightarrow yRx.
\end{align*}
又设 $x,y,z\in G$ 且 $xRy,yRz$, 则 $x\in HyK,y\in HzK$, 于是存在 $h_1,h_2\in H,k_1,k_2\in K$, 使得
\begin{align*}
x=h_1yk_1=h_1h_2zk_1k_2\in HzK\Longrightarrow xRz.
\end{align*}
综上可知 $R$ 是等价关系. 又由
\begin{align*}
gRx\Longleftrightarrow g\in HxK\Longleftrightarrow g=hxk\left( h\in H,k\in K \right) \Longleftrightarrow x=h^{-1}gk^{-1}\in HgK.
\end{align*}
知 $g$ 所在的等价类为 $HgK$.由\refthe{Set Theory-theorem:等价类就和集合的分划对应-定理1.1.1}知 \( \{HgK:g\in G\} \) 为 \( G \) 的一个分划.

{\heiti 必要性:}由等价关系的定义和\refthe{theorem:子群的充要条件--抽象代数}是显然的.
\end{enumerate}
\end{proof}

\begin{theorem}\label{theorem:子群陪集相关性质}
设$A,B,H$是群$G$的非空子群,$a,b\in G$,则
\begin{enumerate}[(1)]
\item $a\in aH$.

\item\label{theorem:子群陪集相关性质-2} $aH=H$的充分必要条件是$a\in H$.

\item $aH$为子群的充分必要条件是$a\in H$.

\item\label{theorem:子群陪集相关性质-4} $aH=bH\iff a^{-1}b\in H\iff b^{-1}a\in H\iff aH\cap bH\neq \varnothing$;
\\
$Ha=Hb\iff ab^{-1}\in H\iff ba^{-1}\in H\iff Ha\cap Hb\neq \varnothing$.

\item\label{theorem:子群陪集相关性质-5} $AB$也是群$G$的子群的充分必要条件是$AB=BA$.
\end{enumerate}
\end{theorem}
\begin{proof}
\begin{enumerate}[(1)]
\item 只需注意到$a=ae\in aH$.

\item 如果$aH=H$,因$a\in aH$,所以$a\in H$.

反之,如果$a\in H$,则$a^{-1}\in H$.从而
\begin{align*}
aH\subseteq H\cdot H=H, \quad
H=(aa^{-1})H=a(a^{-1}H)\subseteq aH,
\end{align*}
所以$aH=H$.

\item 设$aH$为子群.因为$a\in aH$,所以$a^{-1}\in aH$,于是$e=aa^{-1}\in aH$.从而存在$h\in H$,使$e=ah$.所以$a=eh^{-1}=h^{-1}\in H$.

另一方面,如果$a\in H$,则$aH=H$为子群.

\item 如果$aH=bH$,则
\begin{align*}
a^{-1}bH=a^{-1}aH=H,
\end{align*}
从而由\rrefthe{theorem:子群陪集相关性质}{theorem:子群陪集相关性质-2}知,$a^{-1}b\in H$.

反之,如果$a^{-1}b\in H$,则又由\rrefthe{theorem:子群陪集相关性质}{theorem:子群陪集相关性质-2}得$a^{-1}bH=H$,于是
\begin{align*}
aH=a(a^{-1}bH)=(aa^{-1})bH=ebH=bH.
\end{align*}
故$aH=bH\iff a^{-1}b\in H$.交换$a,b$位置即得$bH=aH\iff b^{-1}a\in H$.

若$aH=bH$,则显然有$aH\cap bH\neq \varnothing$.假设$aH\cap bH\neq\varnothing$.任取$g\in aH\cap bH$,则存在$h_1,h_2\in H$,使
\begin{align*}
ah_1=g=bh_2.
\end{align*}
从而
\begin{align*}
aH=a(h_1H)=(ah_1)H=(bh_2)H=b(h_2H)=bH.
\end{align*}
故$aH=bH\iff aH\cap bH\neq \varnothing$.

$Ha=Hb$的等价条件同理可证.

\item {\heiti 必要性:}设$AB$为$G$的子群.
对任意的$ab\in AB$,其中$a\in A,b\in B$,有$(ab)^{-1}\in AB$.因而存在$a_1\in A$,
$b_1\in B$,使$a_1b_1=(ab)^{-1}$.从而
\begin{align*}
ab=(a_1b_1)^{-1}=b_1^{-1}a_1^{-1}\in BA,
\end{align*}
所以
$$AB\subseteq BA.$$

反之,对任意的$ba\in BA$,其中$b\in B,a\in A$,有$(ba)^{-1}=a^{-1}b^{-1}\in AB$.于是
\begin{align*}
ba=(a^{-1}b^{-1})^{-1}\in AB,
\end{align*}
所以
$$BA\subseteq AB.$$

这就证明了$AB=BA$.

{\heiti 充分性:}对任意的$a_1b_1,a_2b_2\in AB$,其中$a_i\in A,b_i\in B\ (i=1,2)$,由于$AB=BA$,因此由\rrefthe{theorem:子群陪集的基本性质}{theorem:子群陪集的基本性质-1}和\rrefthe{theorem:子群陪集的基本性质}{theorem:子群陪集的基本性质-3}有
\begin{align*}
a_1b_1(a_2b_2)^{-1}=a_1b_1b_2^{-1}a_2^{-1}=a_1(b_1b_2^{-1})a_2^{-1}\in ABA=A(BA)=A(AB)=(AA)B=AB,
\end{align*}
由此知$AB$是$G$的子群.
\end{enumerate}
\end{proof}

\begin{definition}
设$H,K$是群$G$的子群.
\begin{enumerate}[(1)]
\item 由\rrefthe{theorem:抽象代数-定理 1.3.2}{theorem:抽象代数-定理 1.3.2-1}定义$G$中的等价关系$R$为
\[
aRb \iff a^{-1}b \in H.
\]
将 \( G \) 对等价关系$R$的商集合, 即以左陪集$gH,\,g\in G$为元素的集合记为 \( G/H=\{gH:g\in G\} \), 称为 \( G \) 对 \( H \) 的\textbf{左陪集空间}. \( G/H \) 中元素个数 \( |G/H| \) 称为 \( H \) (在 \( G \) 中)的\textbf{指数}, 记为 \( [G:H] \). 

这个等价关系$R$的\hyperref[Set Theory-definition:等价类的完全代表系]{完全代表系}就称为关于$H$的\textbf{左陪集代表元系}或\textbf{左陪集代表系}.显然左陪集代表系的阶就是\( [G:H] \). 

\item 由\rrefthe{theorem:抽象代数-定理 1.3.2}{theorem:抽象代数-定理 1.3.2-2}定义$G$中的等价关系$R$为
\[
aRb \iff ab^{-1} \in H.
\]
将 \( G \) 对等价关系$R$的商集合, 即以右陪集$Hg,\,g\in G$为元素的集合记为 \( G/H=\{Hg:g\in G\} \), 称为 \( G \) 对 \( H \) 的\textbf{右陪集空间}. 

这个等价关系$R$的\hyperref[Set Theory-definition:等价类的完全代表系]{完全代表系}就称为关于$H$的\textbf{右陪集代表元系}或\textbf{右陪集代表系}.显然右陪集代表系的阶就是\( [G:H] \). 

\item 由\rrefthe{theorem:抽象代数-定理 1.3.2}{theorem:抽象代数-定理 1.3.2-3}定义$G$中的等价关系$R$为
\[
xRy \iff x\in HyK.
\]
将 \( G \) 对等价关系$R$的商集合, 即以双陪集$HgK,\,g\in G$为元素的集合记为 \( G/R=\{HgK:g\in G\} \), 称为 \( G \) 对 \( H \) 的\textbf{双陪集空间}. 

这个等价关系$R$的\hyperref[Set Theory-definition:等价类的完全代表系]{完全代表系}就称为关于$(H,K)-$\textbf{双陪集代表元系}或\textbf{双陪集代表系}.显然双陪集代表系的阶就是\( |G/R| \). 
\end{enumerate}
\end{definition}
\begin{remark}
\(\{1\}\) 作为 \( G \) 的子群, 在 \( G \) 中指数显然为 \( |G| \). 故也记 \( |G| = [G: 1] \).
\end{remark}

\begin{example}
设 \( V \) 是数域 \( \mathbb{P} \) 上的 \( n \) 维线性空间, \( GL(V) \) 有子群 \( SL(V) \). 在 \( V \) 中取定一组基, 任何一个线性变换由它在这组基下的矩阵完全确定, 可把它们等同起来.\( \forall \lambda \in \mathbb{P}, \lambda \neq 0, \) 令 \( D(\lambda) = \text{diag}(\lambda, 1, \cdots, 1) \), 于是 \( D(\lambda) \in GL(V) \), 对于 \( A \in GL(V) \) 有
\[
A SL(V) = D(\lambda) SL(V)\iff \ \det A = \lambda.
\]
于是
\[
GL(V) = \bigcup_{\lambda \neq 0} D(\lambda) SL(V),
\]
因而
\[
[GL(V): SL(V)] = +\infty.
\]
\end{example}
\begin{proof}

\end{proof}

\begin{example}
设 \( V \) 是 \( n \) 维 Euclid 空间. 由 \( A \in O(V) \) 有 \( \det A = \pm 1 \),令 \( D(\lambda) = \text{diag}(\lambda, 1, \cdots, 1) \),  于是
\[
O(V) = SO(V) \bigcup D(-1) SO(V),
\quad
[O(V): SO(V)] = 2.
\]
\end{example}
\begin{proof}

\end{proof}

\begin{theorem}\label{theorem:抽象代数-定理 1.3.4}
设 \( H \) 是群 \( G \) 的子群, 则 \( G \) 中由
\[
a R b, \ \text{若} \ a^{-1}b \in H
\]
所定义的关系 \( R \) 为同余关系的充分必要条件是
\[
g h g^{-1} \in H, \quad \forall g \in G, h \in H.
\]
此时称 \( H \) 为 \( G \) 的\textbf{正规子群},记为 \( H \lhd  G \). 
同时, 商集合 \( G/H \) 对同余关系 \( R \) 导出的运算
\begin{align*}
aH \cdot bH = abH, \quad \forall a, b \in G
\end{align*}
也构成一个群, 称为 \( G \) 对 \( H \) 的\textbf{商群}.商群 \( G/H \) 的幺元为 \( 1 \cdot H = H \).为方便计, 常将商群 \( G/H \) 中元素记为 \( \overline{g} = gH \).有时也将商群$G/H$记作$\frac{G}{H}$.并且$H$为$G/H$的幺元,$aH$的逆元为$a^{-1}H$.
\end{theorem}
\begin{proof}
设 \( R \) 为同余关系. 又 \( g \in G, h \in H \), 于是有
\[
gRgh, \quad g^{-1}Rg^{-1},
\]
因而 \( gg^{-1}R(ghg^{-1}) \), 即 \( 1  R  ghg^{-1} \), 亦即 \( ghg^{-1} \in H \).

反之, 设 \( \forall g \in G, h \in H \) 有 \( ghg^{-1} \in H \). 设 \( a R b, c R d \), 则$a^{-1}b,c^{-1}d\in H$,即 \( \exists h_1, h_2 \in H \), 使 \( b = a h_1 \), \( d = c h_2 \), 从而$c^{-1}=h_2d^{-1}$.因而
\( (ac)^{-1}(bd)=c^{-1}a^{-1}ah_1d=h_2\left( d^{-1}h_1d \right) \in H\), 则有 \( (ac)  R  (bd) \), 即 \( R \) 为同余关系.

设 \( R \) 为同余关系. 因 \( a \) 所在等价类为 \( aH \), 由\refthe{Set Theory-theorem:同余关系诱导商集中的乘法-定理1.1.3} 知 \( G/H \) 中的乘法为
\begin{align}
aH \cdot bH = abH, \quad \forall a, b \in G. \label{eq:1.3.3}
\end{align}
显然有 \( (aH \cdot bH)cH = abcH = aH(bH \cdot cH) \), \( 1H \cdot aH = aH \), \( a^{-1}H \cdot aH = 1 \cdot H \), 故 \( G/H \) 为群.
\end{proof}

\begin{definition}[极大正规子群]
设$G$是一个群,$H\lhd G$,如果$H$满足以下两个条件:
\begin{enumerate}[(1)]
\item $H\subset G$,即$H\neq G$.

\item 若$K\lhd G$且$H\subseteq K\subseteq G$,则必有$K=H$或$K=G$.
\end{enumerate}
则称$H$是$G$的\textbf{极大正规子群}.
\end{definition}

\begin{theorem}\label{theorem:抽象代数--半群中的同余关系可导出商集合也是商半群}
如果关系 \( \sim \) 是幺半群 (或半群)$G$中的同余关系, 那么商集合$G/\sim$对导出的运算(见\refthe{Set Theory-theorem:同余关系诱导商集中的乘法-定理1.1.3})也是幺半群 (或半群), 称之为\textbf{商幺半群}(或\textbf{商半群}). 

若$G$是交换幺半群(或交换半群),则商集合$G/\sim$对导出的运算也是交换幺半群(或交换半群).
\end{theorem}
\begin{proof}

\end{proof}

\begin{theorem}\label{theorem:抽象代数-定理 1.3.5}
设 \( H \) 是群 \( G \) 的子群, 则下列条件等价:
\begin{enumerate}[(1)]
\item\label{theorem:抽象代数-定理 1.3.5-1} \( H \lhd  G \);

\item\label{theorem:抽象代数-定理 1.3.5-3} 对$\forall a,b\in G$,如果$ab\in H$,则$ba\in H$;

\item\label{theorem:抽象代数-定理 1.3.5-2} $g h g^{-1} \in H,\forall g \in G, h \in H\iff gHg^{-1}\subseteq H,\forall g\in G$;

\item\label{theorem:抽象代数-定理 1.3.5-4} \( gHg^{-1} = H, \forall g \in G \);

\item\label{theorem:抽象代数-定理 1.3.5-5} \( gH = Hg, \forall g \in G \iff GH=HG\);

\item\label{theorem:抽象代数-定理 1.3.5-6} \( g_1Hg_2H = g_1g_2H, \forall g_1, g_2 \in G \).
\end{enumerate}
\end{theorem}
\begin{remark}
由这个\rrefthe{theorem:抽象代数-定理 1.3.5}{theorem:抽象代数-定理 1.3.5-4}可知\textbf{一个群的任意正规子群都是Abel群.}
\end{remark}
\begin{proof}

\ref{theorem:抽象代数-定理 1.3.5-1} $\Rightarrow$ \ref{theorem:抽象代数-定理 1.3.5-3}. 设$H$是$G$的正规子群,则对任意的$a,b\in G$,如果$ab\in H$,则
\begin{align*}
ba = b(ab)b^{-1}\in H.
\end{align*}

\ref{theorem:抽象代数-定理 1.3.5-3} $\Rightarrow$ \ref{theorem:抽象代数-定理 1.3.5-2}.对$\forall a,b\in G$,如果由$ab\in H$,可推出$ba\in H$,则对任意的$a\in G$,
$h\in H$,由于$a^{-1}(ah)=h\in H$,因此
\begin{align*}
aha^{-1} = (ah)a^{-1}\in H,
\end{align*}

\ref{theorem:抽象代数-定理 1.3.5-2} $\Rightarrow$ \ref{theorem:抽象代数-定理 1.3.5-4}.对$\forall g\in G$,由$gHg^{-1}\subseteq H$知,对$\forall h\in H$,有$g^{-1}hg=\left( g^{-1} \right) ^{-1}hg^{-1}\in H.$
从而$h=g\left( g^{-1}hg \right) g^{-1}\in gHg^{-1}.$
故由$h$的任意性知$H\subseteq gHg^{-1}$.因此$gHg^{-1}=H$,$\forall g\in G$.

\ref{theorem:抽象代数-定理 1.3.5-4} $\Rightarrow$ \ref{theorem:抽象代数-定理 1.3.5-5}.  \( \forall g \in G, h \in H \) 有 \( gh = ghg^{-1}g \in Hg, hg = gg^{-1}hg \in gH \), 故 \( gH = Hg \).

\ref{theorem:抽象代数-定理 1.3.5-5} $\Rightarrow$ \ref{theorem:抽象代数-定理 1.3.5-6}. 设 \( g_1, g_2 \in G, h_1, h_2, h \in H \). 由$gH=Hg$知 \( \exists h_1', h' \in H \), 使 \( h_1g_2 = g_2h_1' \), \( g_2h = h'g_2 \). 于是 \( g_1h_1g_2h_2 = g_1g_2h_1'h_2 \in g_1g_2H \), \( g_1g_2h = g_1h'g_2 \cdot 1 \in g_1H \cdot g_2H \), 故 \( g_1H \cdot g_2H = g_1g_2H \).

\ref{theorem:抽象代数-定理 1.3.5-6} $\Rightarrow$ \ref{theorem:抽象代数-定理 1.3.5-1}. 设 \( g \in G, h \in H \), 故有 \( ghg^{-1} \in gHg^{-1}H = gg^{-1}H = H \), 则 \( H \lhd  G \).
\end{proof}

\begin{proposition}\label{proposition:正规子群的基本性质}
\begin{enumerate}[(1)]
\item\label{proposition:正规子群的基本性质-1} Abel 群 \( G \) 的任一子群 \( H \) 都是正规子群, 商群 \( G/H \) 也是 Abel 群.

\item\label{proposition:正规子群的基本性质-2} 若$H$是群$G$的子群且$H\supseteq N,\,N\lhd G$,则$N\lhd H.$

\item\label{proposition:正规子群的基本性质-6} 设$H$和$N$分别是群$G$的子群和正规子群,则$HN=NH$是$G$的子群.

\item\label{proposition:正规子群的基本性质-5} 若$G$是一个群,则$G$的任意正规子群的交$\bigcap_{H\lhd G}{H}\lhd G$.

\item\label{proposition:正规子群的基本性质-3} 设$G$是一个群,且$N_1$, $N_2$, $\cdots$, $N_k\lhd G$,则$N_1 N_2\cdots N_k\lhd G.$

\item\label{proposition:正规子群的基本性质-7} 设$G$是一个群,$N_1,N_2,\cdots N_k\lhd G$,且$N_i\cap \prod_{j=1}^{i-1}{N_j}=\{1\}(i\neq j)$,则对$\forall i,j\in \{1,2,\cdots,k\},$有
\begin{align*}
n_in_j=n_jn_i,\,\,\forall n_i\in N_i,n_j\in N_j.
\end{align*}
并且
\begin{align*}
N_j\subseteq C_G(N_i) \triangleq  \{g\in G\mid gn_i = n_ig ,\forall n_i\in N_i\},\,\,\forall i,j\in \{1,2,\cdots,k\}.
\end{align*}
\end{enumerate}
\end{proposition}
\begin{remark}
正规子群不具有传递性. 例如, $\mathbb{Z}_2 \lhd K_4 \lhd S_4$, 其中 $K_4$ 是 Klein 四元群. 但 $\mathbb{Z}_2$ 不是 $S_4$ 的正规子群.
\end{remark}
\begin{proof}
\begin{enumerate}[(1)]
\item 因为
\begin{align*}
aH = \{ah \mid h \in H\} = \{ha \mid h \in H\} = Ha, \quad \forall a \in G,
\end{align*}
所以$H$是$G$的正规子群.

对$\forall g_1H,g_2H\in G/H$,由\rrefthe{theorem:抽象代数-定理 1.3.5}{theorem:抽象代数-定理 1.3.5-6}有
\begin{align*}
g_1Hg_2H=g_1g_2H=g_2g_1H=g_2Hg_1H.
\end{align*}
故$G/H$也是Abel群.

\item 由\rrefpro{proposition:抽象代数---子群的基本性质}{proposition:抽象代数---子群的基本性质-3}知$N$是$H$的子群.
又由$N \lhd G$知
\begin{align*}
gng^{-1} \in N\subseteq H,\ \forall n \in N,\ g \in H.
\end{align*}
故$N\lhd H.$

\item 由于$N$是$G$的正规子群,因此对任意的$h\in H$,有$hN=Nh$.由此推出,$HN=NH$,从而由\rrefthe{theorem:子群陪集相关性质}{theorem:子群陪集相关性质-5}知,$HN$为$G$的子群.

\item 设$I$为任意指标集,$H_i\lhd G$,$\forall i\in I$.则对$\forall g\in G$,有
\begin{align*}
gH_ig^{-1}\subseteq H_i,\quad \forall i\in I.
\end{align*}
对$\forall h\in \bigcap\limits_{i\in I}{H_i}$,有$h\in H_i$,$\forall i\in I$.从而由上式可得
\begin{align*}
ghg^{-1}\in H_i,\quad \forall i\in I\Longrightarrow ghg^{-1}\in \bigcap\limits_{i\in I}{H_i}.
\end{align*}
进而$g\bigcap\limits_{i\in I}{H_i}g^{-1}\subseteq \bigcap\limits_{i\in I}{H_i}$.故由\rrefthe{theorem:抽象代数-定理 1.3.5}{theorem:抽象代数-定理 1.3.5-3}知$\bigcap\limits_{i\in I}{H_i}\lhd G$.

\item 由$N_i\lhd G,i=1,2,\cdots,k$知
\begin{align*}
gn_ig^{-1}\in N_i,\quad \forall n_i\in N_i,g\in G.
\end{align*}
于是对$\forall n_1n_2\cdots n_k\in N_1N_2\cdots N_k,g\in G$,有
\begin{align*}
g\left(n_1n_2\cdots n_k\right)g^{-1}=\left(gn_1g^{-1}\right)\left(gn_2g^{-1}\right)\cdots\left(gn_kg^{-1}\right)\in N_1N_2\cdots N_k.
\end{align*}
故$N_1N_2\cdots N_k\lhd G.$

\item 设 $n_i \in N_i, n_j \in N_j$, 不妨设 $i < j$, 则
\begin{align*}
n_i n_j n_i^{-1} n_j^{-1} = n_i (n_j n_i^{-1} n_j^{-1}) \in N_i \subseteq N_1 N_2 \cdots N_{j-1};
\end{align*}
\begin{align*}
n_i n_j n_i^{-1} n_j^{-1} = (n_i n_j n_i^{-1}) n_j^{-1} \in N_j,
\end{align*}
所以
\begin{align*}
n_i n_j n_i^{-1} n_j^{-1} \in (N_1 N_2 \cdots N_{j-1}) \cap N_j = \{1\},
\end{align*}
从而 $n_i n_j n_i^{-1} n_j^{-1} = 1$, 由此得 $n_i n_j = n_j n_i$.
进而
\begin{align*}
n_j\in C_G(N_i),
\end{align*}
再由$n_j$的任意性知
\begin{align*}
N_j\subseteq C_G(N_i).
\end{align*}
再由$i,j$的任意性知结论成立.
\end{enumerate}
\end{proof}

\begin{example}
将商群 \( G/H \) 中元素记为 \( \overline{g} = gH \), 则
\begin{enumerate}[(1)]
\item \( SL(V) \lhd  GL(V) \), \( GL(V)/SL(V) = \{\overline{D(\lambda)}|\lambda \neq 0\} \) 且 \( \overline{D(\lambda)}\overline{D(\mu)} = \overline{D(\lambda\mu)} \);

\item \( SO(V) \lhd  O(V) \), \( O(V)/SO(V) = \{\overline{D(1)}, \overline{D(-1)}\} \);

\item \( A_n \lhd  S_n \), \( S_n/A_n = \{\overline{1}, \overline{\sigma}|\sigma \text{ 奇置换} \} \) 且
\[
\overline{1} \cdot \overline{\sigma} = \overline{\sigma} \cdot \overline{1} = \overline{\sigma}, \quad \overline{\sigma} \cdot \overline{\sigma} = \overline{1} \cdot \overline{1} = \overline{1}.
\]

\item 对任意的$a\in G$,由已知条件知,存在$b\in G$,使$aH = Hb$,则$a\in aH = Hb$,从而$a\in Ha\cap Hb$,即$Ha\cap Hb$非空,因此$Ha = Hb$,于是$aH = Ha$.所以由\refthe{theorem:抽象代数-定理 1.3.5}知$H$是$G$的正规子群.
\end{enumerate}
\end{example}
\begin{proof}

\end{proof}

\begin{example}
令 $G$ 是 $n$ 阶有限群, $S$ 是 $G$ 的一个子集, $|S| > \frac{n}{2}$. 试证对任意 $g \in G$, 存在 $a,b \in S$ 使得 $g = ab$.
\end{example}
\begin{proof}
由\refpro{proposition:有限群乘积的阶}知, 对任一 $g \in G$, $gS^{-1}$ 含有 $|S|$ 个元, 这里 $S^{-1}$ 是 $S$ 中元的所有逆元组成的集合. 因为 $|S| > \frac{|G|}{2}$, 故 $gS^{-1} \cap S \neq \varnothing$. 从而存在 $a,b \in S$, 使得 $gb^{-1} = a$, 即 $g = ab$.
\end{proof}

\begin{theorem}[Lagrange定理]\label{theorem:抽象代数--Lagrange定理}
设 \( H \) 是有限群 \( G \) 的子群, 记1为$G,H$的幺元,则有
\begin{align}
[G:1] = [G:H][H:1] \label{eq:1.3.229034890--1}
\end{align}
即
\begin{align*}
[G:H]=\frac{|G|}{|H|}.
\end{align*}
因而子群 \( H \) 的阶是群 \( G \) 的阶的因子,即$|H|\mid |G|$.
\end{theorem}
\begin{remark}
这个结论对无限群 \( G \) 也正确, 此时等式两边都是 \( +\infty \).
\end{remark}
\begin{proof}
设 \( a \in G \). 显然, 映射 \( h \to ah \) 是 \( H \) 到 \( aH \) 上的一一对应, 因而 \( |aH| = |H| = [H:1] \). 又由\refthe{theorem:抽象代数-定理 1.3.2}知 \( G = \bigcup\limits_{a \in G} aH \) 为不相交的并, \( \{aH:a\in G\} \) 的不同左陪集个数为 \( [G:H] \), 故式 \(\eqref{eq:1.3.229034890--1}\) 成立.
\end{proof}

\begin{corollary}\label{corollary:抽象代数-推论 1.3.5}
\begin{enumerate}[(1)]
\item\label{corollary:抽象代数-推论 1.3.5-0} 设$H$是有限群$G$的子群,$K$是$H$的子群,则$[G:K]=[G:H][H:K]$,进而$[G:H]=\frac{[G:K]}{[H:K]}$.

\item\label{corollary:抽象代数-推论 1.3.5-1} 设 \( G \) 为有限群, \( H \lhd  G \), 则商群 \( G/H \) 的阶 \( [G/H : H] = [G : H] = \frac{[G:1]}{[H:1]} \).

\item\label{corollary:抽象代数-推论 1.3.5-2} 设 \( G \) 为无限群, \( H \lhd  G \), 则商群 \( G/H \) 的阶 \( [G/H : H] = [G : H] \).

\item\label{corollary:抽象代数-推论 1.3.5-3} 设$A,B$是群$G$的有限子群,则$|A|\cdot |B| = |AB|\cdot |A\cap B|$.进而$|AB|=\frac{|A||B|}{|A\cap B|}$.

\item\label{corollary:抽象代数-推论 1.3.5-4} 设$G$是一个群,$H$是$G$的子群,且$[G:H]=\frac{|G|}{|H|}=1$,则$G=H.$
\end{enumerate}
\end{corollary}
\begin{proof}
\begin{enumerate}[(1)]
\item 这是\hyperref[theorem:抽象代数--Lagrange定理]{Lagrange定理}的直接推论. 

\item 这是\hyperref[theorem:抽象代数--Lagrange定理]{Lagrange定理}的直接推论. 

\item 这是\hyperref[theorem:抽象代数--Lagrange定理]{Lagrange定理}的直接推论. 

\item 易知$AB = \bigcup\limits_{a \in A} aB$. 记 $\Sigma = \{aB \mid a \in A\}$, 则 $|AB| = |\Sigma| \cdot |B|$. 令
\begin{align*}
&\phi : A/A\cap B\longrightarrow \Sigma 
\\
&x(A\cap B)\longmapsto xB,\quad \forall x\in A.
\end{align*}
\begin{enumerate}[(i)]
\item 如果 $x(A\cap B) = y(A\cap B)(x,y \in A)$, 则 $y^{-1}x \in A\cap B \subseteq B$, 从而 $xB = yB$, 所以 $\phi$ 为 $A/A\cap B$ 到 $\Sigma$ 的良定义的映射.

\item 设 $x(A\cap B),y(A\cap B) \in A/A\cap B$, 有 $\phi(x(A\cap B)) = \phi(y(A\cap B))$, 即 $xB = yB$, 则 $y^{-1}x \in B$. 又因为 $x,y \in A$, 而 $A$ 为群, 所以 $y^{-1}x \in A$, 于是 $y^{-1}x \in A\cap B$. 从而 $x(A\cap B) = y(A\cap B)$. 所以 $\phi$ 为单映射.

\item 对任意的 $x \in A$, 有 $x(A\cap B) \in A/A\cap B$, 且 $\phi(x(A\cap B)) = xB$, 所以 $\phi$ 为满映射.
\end{enumerate}
由此得 $|A/A\cap B| = |\Sigma|$, 所以再利用\hyperref[theorem:抽象代数--Lagrange定理]{Lagrange定理}可得
\begin{align*}
|AB| \cdot |A\cap B| = |B| \cdot |\Sigma| \cdot |A\cap B| = |B| \cdot |A/A\cap B| \cdot |A\cap B| = |A| \cdot |B|.
\end{align*}

\item 由$[G:H]=1$知$|G/H|=1$,又$H\in G/H$,故$G/H=H$,即对$\forall g\in G,$都有$gH = H$,因此$g\in H$,于是$G\subseteq H$.又$H\subseteq G$,故$G=H$.
\end{enumerate}

\end{proof}

\begin{proposition}\label{proposition:左右陪集的基本性质}
\begin{enumerate}[(1)]
\item\label{proposition:左右陪集的基本性质-1} 设$A$和$B$均为群$G$的子群,则
\begin{enumerate}[(i)]
\item\label{proposition:左右陪集的基本性质-1-1} $g(A\cap B)=gA\cap gB$, $\forall\,g\in G$.

\item\label{proposition:左右陪集的基本性质-1-2} 若$A$和$B$均有有限的指数,即$[G:A],[G:B]$有限,则$A\cap B$也有有限的指数,即$[G:A\cap B]$有限.
\end{enumerate}

\item\label{proposition:左右陪集的基本性质-2} 如果$R$是群$G$对于子群$A$的右陪集代表元系,则$R^{-1}$是群$G$对于$A$的左陪集代表元系.

\item\label{proposition:左右陪集的基本性质-3} 设$A< G$, $B<G$.如果存在$a,b\in G$,使得$Aa=Bb$,则$A=B$.

\item\label{proposition:左右陪集的基本性质-4} 设 $H$ 和 $K$ 分别是有限群 $G$ 的两个子群, $g\in G$,则存在$k_1,k_2,\cdots,k_t\in K$和$h_1,h_2,\cdots,h_s\in H$,其中$t=[K:K\cap g^{-1}Hg],s=[H:H\cap gKg^{-1}]$,使得
\begin{gather*}
HgK=\bigsqcup_{1\leqslant i\leqslant t}{Hgk_i},\quad
HgK=\bigsqcup_{1\leqslant i\leqslant s}{h_igK}.
\end{gather*}
并且
\begin{gather*}
|HgK| = |H|[K : K\cap g^{-1}Hg] = |K|[H : H\cap gKg^{-1}].
\end{gather*}

\item\label{proposition:左右陪集的基本性质-5} 设 $A$ 是群 $G$ 的具有有限指数的子群, 则存在 $G$ 的一组元 $g_1,g_2,\cdots,g_m$, 它们既可以作为 $A$ 在 $G$ 中的右陪集代表元系, 又可以作为 $A$ 在 $G$ 中的左陪集代表元系.
\end{enumerate}
\end{proposition}
\begin{proof}
\begin{enumerate}[(1)]
\item \begin{enumerate}[(i)]
\item $g(A\cap B)\subseteq gA$, $g(A\cap B)\subseteq gB$, 故$g(A\cap B)\subseteq gA\cap gB$. 反之,
$\forall\,x=gh\in gA\cap gB$, 则$h\in A\cap B$, 从而$x=gh\in g(A\cap B)$. 于是$g(A\cap B)=gA\cap gB$, $\forall\,g\in G$.

\item {\color{blue}证法一:}\rrrefpro{proposition:左右陪集的基本性质}{proposition:左右陪集的基本性质-1}{proposition:左右陪集的基本性质-1-1}表明$A\cap B$的任一左陪集是$A$的一个左陪集和$B$的一个左陪集之交. 若$A$和$B$均有有限个左陪集,则$A\cap B$也只有有限个左陪集. 由此即得证.

{\color{blue}证法二:}用$\langle AB\rangle$表示$G$的由$AB$生成的子群.则$|\langle AB\rangle|\geqslant |AB|=\frac{|A||B|}{|A\cap B|}$. 两边同乘$\frac{|G|}{|A||B|}$,再利用\hyperref[theorem:抽象代数--Lagrange定理]{Lagrange定理}得到
\begin{align*}
\left[ G:A\cap B \right] =\frac{|G|}{|A\cap B|}\leqslant \frac{|G|}{|A|}\cdot \frac{|\langle AB\rangle |}{|B|}\leqslant \frac{|G|}{|A|}\cdot \frac{|G|}{|B|}<+\infty .
\end{align*}
\end{enumerate}

\item 由题设知$G=\bigcup\limits_{g\in R}Ag$,并且满足$Ag\cap Ah=\varnothing,\forall g\neq h,\,g,h\in R$.需要验证$G=\bigcup\limits_{x\in R^{-1}}xA$,并且满足$xA\cap yA=\varnothing,\,\forall x\neq y,\,x,y\in R^{-1}$.

显然$G\supseteq \bigcup\limits_{x\in R^{-1}}xA$.$\forall z\in G$,由题设$z^{-1}\in Ag$,对某一$g\in R$,于是$z\in g^{-1}A^{-1}=g^{-1}A$.故$G\subseteq \bigcup\limits_{x\in R^{-1}}xA$.因此$G = \bigcup\limits_{x\in R^{-1}}xA$.
若$z\in xA\cap yA$,其中$x,y\in R^{-1}$且$x\neq y$,即$z=xa=yb$, $a,b\in A$,从而
\begin{align*}
Ax^{-1}\cap Ay^{-1}=A(az^{-1})\cap A(bz^{-1})=Az^{-1}\neq\varnothing
\end{align*}
其中$x^{-1},y^{-1}\in R$.从而由题设$x^{-1}=y^{-1}$, 进而$x=y$矛盾!

\item 因$A=Bba^{-1}$,故$ba^{-1}=1\cdot ba^{-1}\in Bba^{-1}=A$,但$A$是子群,故$(ba^{-1})^{-1}=ab^{-1}\in A$.于是$A=Aab^{-1}=Bbb^{-1}=B$.

\item 定义$K$上的一个关系$\sim$:
$$k\sim k'\Longleftrightarrow Hgk=Hgk'.$$
显然这个关系是等价关系.设这个等价关系的完全代表系为$\{ k_1,k_2,\cdots,k_t \}$,则由\refthe{Set Theory-theorem:等价类就和集合的分划对应-定理1.1.1}知
\begin{align}
K=\bigsqcup_{1\leqslant i\leqslant t}{\left[ k_i \right]},\,\,\text{其中}\left[ k_i \right] =\left\{ k'\in K\mid Hgk_i=Hgk' \right\} .\label{eq::::-289huw8f389hw3fuszfbgnmm21.1}
\end{align}
从而$[k_i]\cap [k_j]=\varnothing,i\neq j$.于是$Hgk_i\neq Hgk_j,i\neq j$,故由\rrefthe{theorem:子群陪集相关性质}{theorem:子群陪集相关性质-4}知$Hgk_i\cap Hgk_j = \varnothing,i\neq j$.

对$\forall hgk\in HgK$,有$h\in H,k\in K$,则由\eqref{eq::::-289huw8f389hw3fuszfbgnmm21.1}式知,存在$i\in [1,t]\cap \mathbb{N}$,使得$k\in [k_i]$,即$Hgk=Hgk_i$.从而$hgk\in Hgk=Hgk_i$.故$HgK\subseteq \bigcup\limits_{1\leqslant i\leqslant t}Hgk_i$.又显然$HgK\supseteq \bigcup\limits_{1\leqslant i\leqslant t}Hgk_i$.因此
\begin{align*}
HgK=\bigsqcup\limits_{1\leqslant i\leqslant t}Hgk_i.
\end{align*}
因此由\refpro{proposition:有限群乘积的阶}知
\begin{align*}
|HgK|=\sum\limits_{i=1}^t|Hgk_i|=|H|\cdot t.
\end{align*}
注意到对$\forall k,k'\in K$,有
\begin{align*}
k\sim k'\Longleftrightarrow Hgk=Hgk'\Longleftrightarrow gkk'^{-1}g^{-1}\in H\Longleftrightarrow kk'^{-1}\in K\cap g^{-1}Hg.
\end{align*}
故由\rrefthe{theorem:抽象代数-定理 1.3.2}{theorem:抽象代数-定理 1.3.2-2}知
\begin{align*}
[k]=\{ k'\in K\mid Hgk=Hgk' \}=(K\cap g^{-1}Hg)k,\quad \forall k\in K.
\end{align*}
因此由\eqref{eq::::-289huw8f389hw3fuszfbgnmm21.1}式可得
\begin{align*}
K=\bigsqcup\limits_{1\leqslant i\leqslant t}[k_i]=\bigsqcup\limits_{1\leqslant i\leqslant t}(K\cap g^{-1}Hg)k_i.
\end{align*}
于是$t=[K:K\cap g^{-1}Hg]$.从而由\rrefcor{corollary:抽象代数-推论 1.3.5}{corollary:抽象代数-推论 1.3.5-0}可得
\begin{align*}
|HgK|=|H|\cdot [K:K\cap g^{-1}Hg].
\end{align*}
同理可得
\begin{align*}
HgK=\bigsqcup\limits_{1\leqslant i\leqslant s}h_igK,\,\,\text{其中}s=[H:H\cap gKg^{-1}].
\end{align*}
故由\refpro{proposition:有限群乘积的阶}知
\begin{align*}
|HgK|=\sum_{i=1}^{s}|h_igK|=|K|\cdot[H:H\cap gKg^{-1}].
\end{align*}

\item 由\rrefthe{theorem:抽象代数-定理 1.3.2}{theorem:抽象代数-定理 1.3.2-3},可将 $G$ 写成双陪集的无交并: $G = \bigsqcup_{g\in R} AgA$,其中$R$是$(A,A)$-的双陪集代表系.对$\forall g\in G$,由\rrefpro{proposition:左右陪集的基本性质}{proposition:左右陪集的基本性质-4}知每个双陪集 $AgA$ 既是 $[A:A\cap gAg^{-1}]$ 个互不相交右陪集之并, 也是 $[A:A\cap gAg^{-1}]$ 个互不相交左陪集之并. 记 $t = [A:A\cap gAg^{-1}]$, 则有无交并
\begin{align*}
AgA = \bigsqcup_{1\leqslant i\leqslant t} Aga_i = \bigsqcup_{1\leqslant i\leqslant t} b_igA,
\end{align*}
其中 $a_i, b_i \in A$. 于是 $A = Ab_i = a_iA$, 故有无交并
\begin{align*}
AgA = \bigsqcup_{1\leqslant i\leqslant t} Ab_iga_i,\quad AgA = \bigsqcup_{1\leqslant i\leqslant t} b_iga_iA.
\end{align*}
因为$a_i,b_i,t$与$g$有关,所以记 $b_i = b_i(g), a_i = a_i(g), t = t(g)$, 于是
\begin{align*}
G=\bigsqcup_{g\in R}{AgA}=\bigsqcup_{g\in R}{\bigsqcup_{1\leqslant i\leqslant t\left( g \right)}{Ab_i\left( g \right) ga_i\left( g \right)}}=\bigsqcup_{g\in R}{\bigsqcup_{1\leqslant i\leqslant t\left( g \right)}{b_i\left( g \right) ga_i\left( g \right) A}}.
\end{align*}
因此
\begin{align*}
\bigcup_{g\in R} \{ b_i(g)ga_i(g) \mid 1 \leqslant i \leqslant t(g) \}
\end{align*}
既是 $G$ 关于 $A$ 的右陪集的一个代表元系, 也是 $G$ 关于 $A$ 的左陪集的一个代表元系.
\end{enumerate}
\end{proof}















\end{document}